\documentclass[11pt,a4paper]{article}

\usepackage[T1]{fontenc}
\usepackage[utf8]{inputenc}
\usepackage[spanish]{babel}
\usepackage[a4paper,margin=2.8cm]{geometry}
\usepackage{microtype}
\usepackage{xurl}
\usepackage[hidelinks]{hyperref}

\hypersetup{
  pdftitle={Paraíso perdido},
  pdfauthor={Juan José Gómez Cadenas}
}

\setlength{\parindent}{1.25em}
\setlength{\parskip}{0.35em}
\raggedbottom

\title{Paraíso perdido}
\author{Juan José Gómez Cadenas}
\date{}

\begin{document}

\maketitle


Hace ahora poco más de un siglo, David Hilbert acuñaba la frase: «Nadie podrá expulsarnos del paraíso que Cantor ha creado para nosotros». Estamos en Münster, durante la conferencia «\emph{Über das Unendliche}»; la fecha es 1925 y Hilbert lleva un cuarto de siglo embarcado en una de las empresas intelectuales más ambiciosas del \emph{novecento}. Quiere rescatar a las matemáticas de las paradojas de la teoría de conjuntos, formuladas por Bertrand Russell y otros entre finales del siglo XIX y principios del XX. Buena parte de esas paradojas tienen su origen en el trabajo de Cantor, quien había descubierto la existencia de un extraño Paraíso, poblado por jerarquías de infinitos ---los llamados cardinales transfinitos--- que los bienaventurados podían manipular. Hilbert tiene como objetivo eliminar las paradojas que amenazan con expulsar a los matemáticos de esa Tierra Prometida. Para ello, propone formalizar por completo las matemáticas y demostrar que un sistema así construido no puede producir contradicciones. Su programa implica desarrollar una metodología que permita demostrar teoremas sin hacer otra cosa, en última instancia, que manipular símbolos. 

Es decir, Hilbert quiere construir una máquina de fabricar teoremas.

Unos pocos años más tarde, Alan Turing reformula la idea en términos de un ingenio capaz de leer instrucciones codificadas en una cinta sin fin ---papel perforado en su época--- y escribir sus resultados en otra cinta análoga. Ese artefacto puede codificar cualquier algoritmo y, por tanto, ejecutar las secuencias de manipulaciones simbólicas que llevan a la demostración de cualquier teorema. De esta manera, la entelequia de Hilbert se transforma en un objeto todavía ideal pero mucho más cerca de lo real. La máquina de Turing. 

Permíteme, incisiva lectora, atento lector, un pequeño \emph{cliffhanger} antes de continuar con esta historia. Saltemos a la noche del 24 de noviembre de 1960, a una cena que podemos imaginar desenfadada y solemne a partes iguales. No faltan ni el \emph{confit} de pato ni una selección de quesos bretones, regada con un buen Borgoña y rematada por un excelente \emph{brandy} y cigarros cubanos. Los anfitriones son el célebre literato y matemático aficionado Raymond Queneau y el ingeniero y ajedrecista François Le Lionnais. Hay una decena de comensales, que incluyen la flor y nata de la \emph{intelligentsia divine} francesa. Todos tienen ganas de divertirse y de liarla un poco, sobre todo después de los calvados. En algún momento de la larga velada, acuerdan formar el Seminario de literatura experimental, que acaba por convertirse en el Oulipo (\emph{Ouvroir de littérature potentielle}, o «taller de literatura potencial»). Los miembros fundadores se comprometen a serlo \emph{ad infinitum}, incluyendo continuidad \emph{post mortem} ---se acuerda que los que pasen a mejor vida podrán eludir participar en futuras reuniones excusándose por defunción---. La broma podría haberse quedado ahí, pero el taller en cuestión se hace famoso y acaba por enrolar a luminarias como Georges Perec, Italo Calvino y Marcel Duchamp.

Los miembros del Oulipo postulan que la literatura, como las matemáticas, puede, en último término, reducirse a un algoritmo capaz de manejar símbolos ---los del lenguaje--- a partir de una serie de reglas fijas. La idea recuerda mucho los fundamentos del programa de Hilbert y el propio Queneau la desarrolla en uno de sus textos, titulado nada menos que «\emph{Les fondements de la littérature d'après David Hilbert}». Queneau toma los axiomas de los \emph{Grundlagen der Geometrie} y sustituye puntos, rectas y planos por palabras, frases y párrafos. Es decir, literalmente, construye teoremas lingüístico-literarios.

Permíteme, curiosa lectora, siempre alerta lector, un poco de contexto. Queneau y los oulipianos se oponen frontalmente al movimiento surrealista de Breton del que el propio Queneau ha formado parte en su juventud. El Oulipo puede explicarse, al menos parcialmente, como la reacción inmunitaria de parte de la élite literaria francesa al reinado ---que ya se iba convirtiendo en tiranía--- de Breton y su corte. Así, donde los unos predican la necesidad de la escritura automática como método para explorar el subconsciente y quizás aprehender lo inefable, los otros se burlan de una técnica muy proclive a producir lo que hoy no dudaríamos en calificar de \emph{slop} y proponen a cambio encadenar al literato a un conjunto de reglas tan estricto que la creatividad no tenga otro remedio que surgir para producir algo respetándolas\footnote{Una propuesta que evoca el célebre chiste en el que un explorador se pasea por la jungla acarreando un pesado yunque, de tal forma que, en caso de encontrarse con un león, puede tirarlo y correr más rápido.}.

En 1961, Queneau publica sus célebres
\emph{Cent mille milliards de poèmes}. Se trata de diez sonetos con las mismas rimas, en los que cada verso se escribe en una tira de papel ofreciendo al lector $10^{14}$ posibles combinaciones. Es, como bien señala mi querido Ángel Fernández en su reciente artículo «Las tijeras de Queneau en la Biblioteca de Babel», una deliciosa vuelta de tuerca al argumento del argentino. Donde Borges solo puede mostrarnos horror combinatorio, Queneau nos ofrece un número gigantesco de posibles poemas, todos los cuales tendrían, al menos formalmente, sentido.

Queneau acaba de construir algo similar a la máquina de Turing. Una máquina de hacer sonetos, capaz de producir un número astronómico de ellos siguiendo un simple algoritmo combinatorio. Por ejemplo: 


\begin{verse}
\itshape
Il se penche il voudrait attraper sa valise\\*
pour la mettre à sécher aux cornes des taureaux\\*
une toge il portait qui n’était pas de mise\\*
il chantait tout de même oui mais il chanta faux
\end{verse}

Traducido al español:

\begin{verse}
Se inclina, quisiera alcanzar su maleta\\*
para ponerla a secar en los cuernos de los toros;\\*
llevaba una toga que no venía al caso,\\*
cantaba de todos modos, sí, pero desafinaba.
\end{verse}

¿Tienen sentido los versos o solo lo parece? Los seres humanos le buscamos sentido a todo y más todavía cuando se trata del lenguaje, así que nos esforzamos por imaginar a un señor vestido con una toga algo desaliñada, cantando desafinadamente mientras se esfuerza en alcanzar su maleta (mojada) para ponerla a secar nada menos que en los cuernos de los toros. Por una parte, el cuarteto no dice nada. Por otra parte, es una pequeña descarga eléctrica. Y nada obliga al lector a escoger los versos al azar. Podemos imaginar que este se torna escritor mediante el acto de seleccionar cuidadosamente los versos que le permitan expresarse, o inquietarnos, o explorar extraños nuevos territorios, que pueden incluir maletas que cuelgan en un enorme tendedero hecho de asta de miuras. 

Sostiene Ángel Fernández que la escritura híbrida entre hombre y máquina es posible, utilizando un sistema similar al de los cien billones de poemas y la propuesta, a primera vista, es muy sugestiva. Uno puede imaginar al LLM como la máquina de producir versos en tiras de papel ---una máquina de Queneau--- y al escritor como aquel que los selecciona. 

Pero creo que, atractiva y todo, la analogía no se sostiene. Y para apoyar mis argumentos, quiero referirme a ese bello juego-deporte que tanto me apasiona ---y en el que tan torpe soy---. Es bien sabido que los sistemas IA que los ajedrecistas llamamos «motores» ---un término que denota que no los confundimos, como solemos hacer con Claude, con nuestro mejor amigo--- nos superan desde hace mucho. Pues bien, en muchas ocasiones se han organizado torneos híbridos, en los que se enfrentaban dos oponentes humanos ---normalmente grandes maestros--- «asistidos» por dos motores de ajedrez. Hubo un tiempo en el que el \emph{hype} en torno a esos torneos era considerable, sobre todo cuando involucraban a grandes jugadores, pero lo cierto es que, hoy en día, no le interesan demasiado a nadie. La razón es muy sencilla. El bando que juega con AlphaZero gana, aunque en el bando del motor inferior juegue el mismísimo Magnus y en el de AlphaZero un pichón como el que suscribe. La moraleja es que la colaboración entre motores de ajedrez y humanos es tan imposible o ridícula como pretender que colaboren Queneau y un niño de seis años en escribir un soneto. 

Magnus Carlsen tiene un Elo que ronda los 2800. AlphaZero anda por los 3500. AlphaZero puede entrenar a Magnus, pero Magnus no aporta nada a AlphaZero.

Pues bien, si Claude tiene un Elo de 2000 como escritor ---y estoy siendo generoso---, Queneau es un 2600, Borges un 2800 y Rainer Maria Rilke pasa de 3000. Todos estos escritores ---y muchos otros de su calibre--- pueden contribuir a entrenar a Claude, pero Claude no tiene nada que aportarles a ellos. 

Claro que ---y ese es el problema--- el zoológico de plumíferos está lleno de pichones que no llegan a los 1000 de Elo y esos, naturalmente, están muy interesados en los concursos híbridos en los que se pueden apoyar en un jugador de 2000... Al menos los más honestos, porque además hay ejércitos de desaprensivos perfectamente dispuestos a hacer trampa y usar el LLM sin darle crédito, \emph{simplemente para impostar un Elo mayor como escritores}. 

Y aquí el ajedrez viene a cuento de nuevo. No hay cosa más absurda que jugar una partida \emph{online} ayudándose de un motor
%La gracia del ajedrez es, precisamente, jugar. Cuando un tramposo deja que el robot mueva por él las piezas, le está haciendo perder el tiempo a su oponente y se está engañando a sí mismo en la misma medida que se engañaría un halterófilo que hiciera sentadillas con 500 kilos..., ayudado por una máquina que se los levantara. 
y, aun así, el ajedrez en la red está infectado por los tramposos. Unos lo hacen por vanidad, en otros casos se trata de impostores ---\emph{influencers} simpáticos que se han hecho un nombre en el negocio floreciente de las sesenta y cuatro casillas a base de montar un buen \emph{show}, pero no saben distinguir un caballo de un alfil--- que defienden un negocio. Lo mismo puede decirse del mundo de la pluma y las plumas. Hay una legión de narcisistas que han descubierto que pueden darle rienda a su \emph{chatbot} y firmar un artículo que da el pego. Hay otra legión de periodistas, articulistas y otras gentes de mal vivir ---y lo de mal vivir es literal, porque se cobra muy poco en esa profesión--- que queman \emph{tokens} como un fumador empedernido quema pitillos porque se ganan el pan juntando palabras y Claude les permite publicar diez artículos con menos esfuerzo del que necesitaban ayer para escribir uno. 

En los clubs de ajedrez se inventó la policía \emph{online} mucho antes de que Substack decidiera utilizar Pangram y por las mismas razones. Para intentar contener a las hordas de narcisos y las avalanchas de \emph{bots}. Para pedir un poco de honestidad al que firma un artículo y al que juega una partida. Como ya he escrito varias veces, no tengo problema ni en jugar con motores ---lo hago a menudo, en mi intento desesperado de mejorar mi ajedrez--- ni en leer texto generado por máquinas, si me parece interesante. Pero sí tengo un gran problema cuando juego con un tramposo que se ayuda con motores y tengo el mismo problema cuando detecto que tras la firma de un humano hay un bot. 

Porque, además, la metáfora que Ángel propone tiene sus limitaciones. Queneau y compañía no delegaban en absoluto la construcción de los versos en otro escritor. Por el contrario, trabajaban denodadamente para construir las estructuras que respetaran las arduas reglas que ellos mismos se habían impuesto y jamás se les habría pasado por la cabeza encargarle ese duro trabajo a otro ---por muy compadres que fueran--- al igual que Breton nunca le pidió a Queneau que utilizara la escritura automática para indagar sobre la psique... de Breton. A los oulipianos les daba casi igual qué se hiciera luego con los versos de sus tiras de papel, ya que para ellos la literatura consistía en construir la máquina, no en darle al botón y hacerle escupir teoremas. 

El escritor que pretende escribir proporcionando un \emph{prompt} al motor y editando el resultado está delegando la construcción de la máquina de Queneau en el LLM y limitándose a hacer el trabajo que podríamos encargar a cualquiera, incluyendo un chimpancé aplicado. 


«El \emph{prompt} captura mis ideas y la máquina me ayuda a escribirlas» es una falacia cósmica. Las ideas no se formulan de un golpe, ni se condensan en un \emph{prompt}. La primera vez que nos cruzan por la cabeza son poco más que una intuición, una vaga pista de algo que no toma forma hasta que lo escribimos. El primer borrador suele ser un batiburrillo de confusas ideas que solo sirve para crear un segundo borrador. Quizás, a partir del tercero ---este artículo ya va por su quinta versión--- empezamos a entender lo que decimos. Delegar un \emph{prompt} a una máquina, digamos: «escribe un artículo que sea un espejo de “Borges en su biblioteca”, utilizando como contraargumento el Oulipo» y luego repasar lo que la máquina produce puede resultar en un texto simpático, provocador y exquisitamente trivial. 

Pero volvamos a Hilbert. Unos años después de su célebre declaración exigiendo el Paraíso de Cantor, Gödel publica sus teoremas de incompletitud que pueden resumirse así: «todo sistema formal consistente contiene enunciados verdaderos que no pueden demostrarse». Poco más tarde, el mismo Turing da una versión informática con su célebre \emph{halting problem} (problema de la parada): no existe un algoritmo capaz de decidir si un programa cualquiera se detendrá alguna vez. Hay preguntas perfectamente bien planteadas que ninguna máquina puede responder.

¿Debemos entender que Gödel y Turing expulsaron a los matemáticos del Paraíso? Yo creo que no. Creo, con Penrose, que los teoremas de incompletitud revelan que la realidad ---y por tanto la mente humana, ya que la realidad pasa por ella--- contiene una parte irreduciblemente no algorítmica. Y en esa irreducibilidad fuera del alcance de un algoritmo ---y de una máquina--- está nuestra última Tierra Prometida.

Ese Paraíso, sin embargo, es un Paraíso perdido. La literatura existe, precisamente, porque somos unos exiliados del jardín del Edén, porque nuestras vidas son fugaces, nuestros destinos inciertos y nuestras penas muy grandes. Existe, como ya he dicho en otra ocasión, para cantar a las cosas que merecen lágrimas. 

Merecen lágrimas la muerte de Dido y la de Julieta, la escena en la que la Maga descubre que su bebé ha dejado de respirar, el páramo que atraviesa un personaje de \emph{Cien años de soledad} donde el eco repite los pensamientos, Príamo besando las manos de Aquiles, el momento en el que Orfeo mira hacia atrás y pierde para siempre a Eurídice, los versos de Dylan Thomas, «\emph{Rage, rage against the dying of the light}», la tenacidad del escribiente que decide no hacerlo, la locura de Ahab, el amor desesperado del padre que lleva una sola bala en su revólver para matar a su hijo antes que abandonarlo a las alimañas que pueblan la Tierra.

Defiendo que ninguna IA ---no las que conocemos--- escribirá líneas comparables a esas. Es posible que sí consigan copiar a los oulipianos, porque los LLM son algoritmos y una escuela que quiere basar la literatura en algoritmos puede acabar recibiendo una dosis y media de ellos. Quizás Claude mejorará un día los cien billones de poemas. Pero sostengo que no captará jamás la esencia última de la literatura, porque, a diferencia nuestra, Claude, bendita sea su alma algorítmica, no ha sido nunca expulsado del Paraíso. 


\end{document}
